\scanpage{64}
\mspage{\handwritten{- 33 -}}
l'amplitude de $z$ comprise entre $0$ et $2\handwritten{\pi}$, $B$ est l'ensemble des $z\handwritten{\neq}0$ tels que $\handwritten{\frac{\pi}{3}<\theta'<\frac{5\pi}{3}}$. L'ensemble ouvert $A\handwritten{\cap}B$ est non connex ses deux composantes connexes $H$ et $H'$ étant respectivement défini par $\handwritten{-\frac{2\pi}{3}<\theta<-\frac{\pi}{3}}$ et $\handwritten{\frac{\pi}{3}<\theta<\frac{2\pi}{3}}$. Soit $f(z)=\handwritten{\sqrt{|z|}}\,e^{i\handwritten{\theta}/2}$ dans $A$, $g(z)=\handwritten{\sqrt{|z|}}\,e^{i(\handwritten{\theta'}-2\handwritten{\pi})/2}$ dans $B$ ; on a $(f(z))^{\handwritten{2}}=z$ dans $A$ et $(g(z))^{\handwritten{2}}=z$ dans $B$, et par suite, comme $f$ et $g$ sont continues, ce sont des fonctions holomorphes dans $A$ et $B$ respectivement (chap.I, \handwritten{\S}\,2, th.1). Mais dans $H$, on a $\handwritten{\theta'}=2\handwritten{\pi}+\handwritten{\theta}$, d'où $f(z)=g(z)$ et dans $H$, $\handwritten{\theta'=\theta}$, d'où $f(z)=-g(z)$.

\handwritten{2)} Lorsque $K=\handwritten{C}$, au lieu de supposer que $\handwritten{f(x)=g(x)}$ dans un ens% right edge clips the rest of this line; no em survives
ble ouvert non vide $U\handwritten{\subset}A\handwritten{\cap}B$, il suffit de supposer que $\handwritten{f(x)=g(x)}$ dans l'intersection d'un tel ensemble $U$ et d'une variét\unsure{é} linéaire réelle $V$ définie par $p$ équations de la forme $\handwritten{\unsure{\mathfrak{I}|\mathscr{I}|\mathscr{J}}}(z_k)=\handwritten{\beta}_k$ $(1\handwritten{\leq}k\handwritten{\leq}p)$, $U\handwritten{\cap}V$ étant supposé non vide. En effet, si $\handwritten{a\in}U\handwritten{\cap}V$, les dérivées partielles de $\handwritten{h}$ au point $a$ sont égales respectivement aux dérivées partielles de mêmes indices de la restriction de $\handwritten{h}$ à $U\handwritten{\cap}V$ ; comme par hypothèse ces dernières sont toutes nulles, la série entière (par rapport aux $x_i-a_i$) égale à $\handwritten{h(x)}$ dans un polycylindre assez petit de centre $\handwritten{a}$, a tous ses coefficients nuls, ce qui montre que $\handwritten{f(x)=g(x)}$ dans un voisina\handwritten{-}ge de $\handwritten{a}$ par rapport à $A\handwritten{\cap}B$.

Si $K=\handwritten{R}$ ou $K=\handwritten{C}$, et $p=1$, il suffit de supposer que $\handwritten{f(x)=g(x\unsure{)}}$ dans une partie $M$ de $A\handwritten{\cap}B$ telle que l'intersection de $M$ et d'un vois\unsure{i}na\handwritten{-}ge compact $U\handwritten{\subset}A\handwritten{\cap}B$ d'un point de $A\handwritten{\cap}B$ soit un ensemble \uline{infini}. En effet, il existe alors un point $\handwritten{b\in}U$ adhérent à une suite de points distincts $(z_n)$ de $U$ tels que $\handwritten{h}(z_n)=0$ pour tout $n$ ; d'après le principe des zéros isolés, on a $\handwritten{h(x)}=0$ dans un voisinage de $\handwritten{b}$ contenu dans $A\handwritten{\cap}B$, et on conclut comme précédemment.
% archive sheet number: 63
